\documentclass[12pt,onecolumn,a4paper]{article}
\usepackage{xcolor}
\usepackage[a4paper,left=20mm,right=20mm,top=25mm,bottom=25mm,headheight=15pt]{geometry}
\usepackage{amsmath,amssymb,mathtools,booktabs,longtable,tabularx,array}
\usepackage{graphicx,float,caption,enumitem,fancyhdr,listings}
\usepackage[most]{tcolorbox}
\usepackage{hyperref}
\usepackage{xepersian}
\settextfont{XB Niloofar}
\setlatintextfont{Arial}
\setmonofont{Courier New}
\providecommand{\UseMathForPositioningText}{}
\setlength{\parindent}{0pt}
\setlength{\parskip}{5pt}
\linespread{1.25}
\hypersetup{colorlinks=true,linkcolor=blue!55!black,urlcolor=blue!55!black,citecolor=magenta!65!black}
\definecolor{codebg}{RGB}{246,246,242}
\definecolor{codeframe}{RGB}{180,180,180}
\lstdefinestyle{project}{backgroundcolor=\color{codebg},frame=single,rulecolor=\color{codeframe},basicstyle=\ttfamily\footnotesize,breaklines=true,showstringspaces=false,numbers=left,numberstyle=\tiny\color{gray},numbersep=6pt,tabsize=2}
\lstset{style=project}

% Custom boxes using tcolorbox
\newtcolorbox{takeawaybox}[1][]{colback=green!3,colframe=green!55!black,title=نتیجه کلیدی (Key Takeaway),fonttitle=\bfseries,breakable,#1}
\newtcolorbox{invariantbox}[1][]{colback=blue!3,colframe=blue!55!black,title=اصل الگوریتمی (Algorithmic Invariant),fonttitle=\bfseries,breakable,#1}
\newtcolorbox{axiombox}[1][]{colback=red!3,colframe=red!55!black,title=بیانیه تحلیل خطا (Error Analysis Axiom),fonttitle=\bfseries,breakable,#1}

\pagestyle{fancy}
\fancyhf{}
\fancyhead[R]{پروژه سوم داده‌کاوی -- گزارش فنی جامع}
\fancyhead[L]{موتور تشخیص سرقت ادبی اسناد دو زبانه}
\fancyfoot[C]{\thepage}
\renewcommand{\headrulewidth}{0.5pt}

\begin{document}
\begin{titlepage}
\centering
\vspace*{1.5cm}
{\Large دانشگاه صنعتی امیرکبیر (پلی‌تکنیک تهران)\par}
\vspace{0.5cm}
{\large دانشکده ریاضی و علوم کامپیوتر\par}
\vspace{1.8cm}
{\Huge\bfseries گزارش فنی پروژه سوم داده‌کاوی\par}
\vspace{0.5cm}
{\LARGE\bfseries موتور تشخیص سرقت ادبی و اسناد مشابه\par}
\vspace{0.3cm}
{\Large\lr{Semantic Duplicate \& Near-Plagiarism Detection Engine}\par}
\vfill
\begin{tabular}{rl}
مدرس: & دکتر فاطمه شاکری\\[4pt]
نام دانشجو: & \dotfill\\[4pt]
شماره دانشجویی: & \dotfill\\[4pt]
نیم‌سال: & بهار ۱۴۰۵
\end{tabular}
\vfill
{\large نسخه ۳٫۱ -- مستندات نهایی، بهینه‌سازی پارامترها و گزارش تحلیل خطا\par}
\end{titlepage}

\tableofcontents
\newpage

\section{چکیده}
در این گزارش فنی، معماری و جزئیات پیاده‌سازی یک موتور تشخیص تشابه اسناد و سرقت علمی دو زبانه (فارسی و انگلیسی) تشریح شده است. سیستم با پیاده‌سازی دو مسیر مستقل \lr{MinHash+LSH+Jaccard} و \lr{TF-IDF Weighted SimHash} از پایه و بدون تکیه بر کتابخانه‌های آماده طراحی شده است. برای ارزیابی عملکرد و بهینه‌سازی پارامترها، یک تحلیل گسترده روی ۵۰۰۰ جفت‌سوال از دیتاست استاندارد \lr{Quora Question Pairs} انجام شد. در فاز تمیزکاری متون، با تغییر سیاست حذف ایست‌واژه‌ها به فرم محافظه‌کارانه جهت حفظ کلمات پرسشی و نفی، تضادهای معنایی شناسایی شدند. بهینه‌سازی پارامترها منجر به کشف پیکربندی ایده آل (\lr{num\_perm=128}، \lr{bands=64}، آستانه جاکارد $0.15$ و حد فاصله همینگ $18$) شد که دقت را به ٪۵۵٫۲۹ و بازیابی را به ٪۵۵٫۵۲ رساند. در نهایت، با افزودن سه ابزار خلاقانه (بهینه‌ساز باند LSH، بن‌واژه‌ساز فارسی سبک و امضای ترکیبی کاراکتری)، میزان بازیابی تا ٪۶۲٫۴۸ افزایش یافت.

\section{اهداف پروژه و معماری یکپارچه}
هدف اصلی این سامانه، شناسایی اسنادی است که دست‌خوش بازنویسی‌های معنایی یا سرقت‌های علمی جزئی شده‌اند. به منظور شکستن سد مرتبه پیچیدگی مقایسه‌ها، یک معماری دو مسیره طراحی شد که مسیر اول (لغوی) از طریق \lr{LSH} مقایسه‌ها را فیلتر کرده و مسیر دوم (برداری) با امضای فشرده اثرانگشت، مقایسه بیتی سریعی را در زمان ثابت مقدور می‌سازد.

\subsection{توجیه ریاضی گلوگاه مقایسه جفت‌به‌جفت (All-Pairs Bottleneck)}
برای یافتن تمامی اسناد مشابه در یک مجموعه به اندازه $n$، مقایسه مستقیم جفت‌به‌جفت اسناد بر اساس شباهت جاکارد (\lr{Jaccard Similarity}) با چالش جدی مقیاس‌پذیری روبرو است. از نظر ریاضی، برای مقایسه همه‌جانبه اسناد، باید تمامی جفت‌های غیرتکراری و مستقل از ترتیب را ارزیابی کنیم. تعداد کل مقایسه‌های مورد نیاز برابر است با انتخاب ۲ از $n$:
\[
\binom{n}{2} = \frac{n(n-1)}{2}
\]
با گسترش اندازه ورودی‌ها و میل کردن $n$ به مقادیر بسیار بزرگ، رابطه فوق به صورت زیر باز می‌شود:
\[
\frac{n(n-1)}{2} = \frac{1}{2}n^2 - \frac{1}{2}n
\]
که در تحلیل رفتار مجانبی و با چشم‌پوشی از جملات مرتبه پایین‌تر و ضرایب ثابت، پیچیدگی محاسباتی آن معادل $\mathcal{O}(n^2)$ خواهد بود.

این رابطه درجه دوم بدان معناست که افزایش خطی اندازه مجموعه داده، زمان اجرای محاسبات را به صورت نمایی و مرتبه دوم افزایش می‌دهد. برای نمونه، در یک مجموعه داده در ابعاد واقعی مانند نسخه کامل دیتاست \lr{Quora Question Pairs} که شامل بیش از $400,000$ سوال است، تعداد کل مقایسه‌ها برابر خواهد بود با:
\[
\binom{400,000}{2} = \frac{400,000 \times 399,999}{2} \approx 80,000,000,000
\]
یعنی حدود ۸۰ میلیارد مقایسه جفت‌به‌جفت مستقل! اگر فرض کنیم هر محاسبه جاکارد لغوی روی متون معمولی تنها ۱ میلی‌ثانیه زمان ببرد، اجرای کامل محاسبات بر روی یک پردازنده تک‌هسته‌ای به بیش از ۹۲۵ روز زمان مداوم نیاز دارد. این گلوگاه شدید محاسباتی ثابت می‌کند که محاسبه مستقیم جاکارد برای مجموعه‌های بزرگ عملاً غیرممکن و غیرقابل اجرا است؛ از این رو استفاده از روش‌های تقریبی کارآمد مانند \lr{MinHash} و فیلترهای \lr{LSH} جهت تبدیل مقایسه‌ها به زمان خطی یا زیرخطی الزامی است.

\section{متدولوژی پیش‌پردازش متن و تمیزکاری اطلاعاتی}

تمیزکاری متن سنگ بنای اصلی در خروجی الگوریتم‌های تشابه است. متون ورودی در کلاس \lr{TextPreprocessor} واقع در فایل \texttt{preprocessing.py}، از مراحل زیر عبور می‌کنند:
\begin{enumerate}
\item \textbf{نرمال‌سازی یونیکد}: تبدیل به قالب استاندارد \lr{NFKC} به منظور یکی کردن کدهای یکسان کاراکترها.
\item \textbf{تبدیل حروف عربی}: نگاشت کاراکترهای عربی متداول مانند «ي»، «ك» به معادل‌های استاندارد فارسی «ی» و «ک».
\item \textbf{تمیزکاری نیم‌فاصله‌ها (\lr{ZWNJ})}: نویسه‌های کنترلی نیم‌فاصله (`\textbackslash u200c`) به فاصله عادی تبدیل می‌شوند تا عبارات و واژه‌ها به درستی تفکیک و پردازش شوند.
\item \textbf{فیلترینگ محافظه‌کارانه ایست‌واژه‌ها}: در داده‌کاوی متون کوتاه (مانند سوالات Quora)، حذف ایست‌واژه‌های پیش‌فرض بسیار مخرب است. به عنوان مثال، حذف واژه‌های پرسشی کلیدی (مانند \lr{what}, \lr{how}, \lr{why}) و کلمات نفی (مانند \lr{not}) معنای جملات را یکسان می‌کند در حالی که منطق آن‌ها متضاد است. برای جلوگیری از این مشکل، لیست ایست‌واژه‌ها بازنویسی و بسیار کوچک گردید تا این کلمات کلیدی حفظ شوند.
\end{enumerate}

\begin{axiombox}
حفظ کلمات نفی و کلمات پرسشی در فاز پیش‌پردازش متون کوتاه، به عنوان یک اصل در سرقت علمی معنایی، مانع از هم‌گرایی صوری ساختار جملات متضاد و کاهش مثبت‌های کاذب (FP) می‌گردد.
\end{axiombox}

\section{تحلیل ریاضی و پیاده‌سازی مسیر اول: MinHash و LSH}

\subsection{شینگل‌سازی کلمه‌ای}
متون نرمال‌شده به شینگل‌های ۳-تایی کلماتی متوالی شکسته می‌شوند. برای متون بسیار کوتاه که تعداد توکن‌های آنها کمتر از ۳ است، از سیاست شینگل‌سازی کوتاه استفاده می‌شود که در آن کل متون به عنوان یک شینگل منفرد استخراج می‌گردد تا از از دست رفتن اطلاعات جلوگیری شود.

\subsection{امضای MinHash}
برای تبدیل مجموعه‌های شینگل بزرگ به بردارهای امضای متراکم، از ۲۵۶ تابع هش مستقل از خانواده زیر استفاده می‌شود:
\[
h_i(x) = (a_i x + b_i) \bmod p
\]
که در آن $p = 2^{61}-1$ (یک عدد اول مرسن بسیار بزرگ) است تا احتمال هرگونه تداخل و برخورد هش‌ها به صفر میل کند. ضرایب $a_i$ و $b_i$ به صورت تصادفی با بذر ثابت ۴۲ تولید می‌شوند. الگوریتم هشر پایدار \lr{BLAKE2b} به کار رفته در سیستم، از نوسان مقادیر به علت بذر سیستم‌عامل جلوگیری کرده و بازتولیدپذیری ۱۰۰ درصدی را تضمین می‌کند.

\subsection{بندبندی LSH و اثبات کاهش محاسبات}
امضای متراکم به $b$ باند تقسیم می‌شود که هر باند شامل $r$ سطر است. دو سند زمانی به عنوان کاندید انتخاب می‌شوند که حداقل در یک باند هم‌پوشانی کامل داشته باشند. احتمال کاندید شدن به عنوان تابعی از شباهت واقعی جاکارد $s$ برابر است با:
\[
P(\text{candidate}) = 1 - (1 - s^r)^b
\]
در پیکربندی بهینه انتخاب‌شده (\lr{Configuration A}) با $b=64$ و $r=2$، احتمال کاندید شدن برای جفت‌هایی با تشابه $0.15$ برابر است با:
\[
P = 1 - (1 - 0.15^2)^{64} = 1 - 0.9775^{64} \approx 0.767\quad (76.7\%)
\]
این پیکربندی تضمین می‌کند که نرخ بازیابی مدل در فاز کاندیدسازی بالا بماند؛ سپس تصمیم نهایی با آستانه جاکارد دقیق $0.15$ اتخاذ می‌شود. از نظر تئوریک، این فرآیند محاسبات جفت‌به‌جفت دقیق را از $\mathcal{O}(n^2)$ به تعداد جفت‌های کاندید کاهش می‌دهد که برای مجموعه داده ۵۰۰۰ تایی منجر به کاهش \textbf{٪۷۰٫۶۲} کل مقایسه‌ها شد.

\begin{invariantbox}
بندبندی LSH به عنوان یک فیلتر پیش‌نیاز عمل می‌کند که بدون محاسبه شباهت دقیق جفت‌به‌جفت اسناد، جفت‌های با شباهت پایین را حذف کرده و پیچیدگی محاسباتی مقایسه را به شدت کاهش می‌دهد.
\end{invariantbox}

\section{تحلیل ریاضی و پیاده‌سازی مسیر دوم: TF-IDF Weighted SimHash}

در این مسیر، ویژگی‌ها با ترکیب فرکانس کلمه‌ای زیرخطی و فرکانس معکوس سند هموارسازی شده وزن‌دهی می‌شوند:
\[
\text{TF}(t,d) = 1 + \ln f_{t,d}
\]
\[
\text{IDF}(t) = \ln\left(\frac{N+1}{\text{df}(t)+1}\right) + 1
\]
برداری با طول ۶۴ (متناظر با ۶۴ بیت اثرانگشت) به صورت صفر مقداردهی می‌شود. برای هر ویژگی در سند:
\begin{enumerate}
\item مقدار هش ۶۴ بیتی پایدار آن با \lr{BLAKE2b} محاسبه می‌شود.
\item برای هر بیت $i$ از این هش، در صورت ۱ بودن بیت، مقدار وزن ویژگی به خانه $i$ بردار اضافه شده و در صورت ۰ بودن، از آن کم می‌شود.
\end{enumerate}
اثرانگشت نهایی با بررسی علامت حاصل‌جمع هر بعد ساخته می‌شود:
\[
f_i = 
\begin{cases} 
1 & \text{if } accumulator[i] \ge 0 \\
0 & \text{if } accumulator[i] < 0 
\end{cases}
\]
فاصله همینگ بین دو اثرانگشت با عملیات بیتی \lr{popcount(x XOR y)} محاسبه می‌شود. با کاهش مرز فاصله همینگ از ۲۵ به ۱۸، مثبت‌های کاذب ناشی از کلمات عمومی حذف شده و دقت SimHash به شدت بهبود یافت.

\section{جست‌وجوی شبکه‌ای و ارزیابی عملکرد (Quora Dataset)}

جدول \ref{tab:grid} ارزیابی کامل جست‌وجوی شبکه‌ای روی ۵۰۰۰ جفت‌سوال دیتاست Quora را نشان می‌دهد:

\begin{table}[H]
\centering
\caption{نتایج جست‌وجوی شبکه‌ای پارامترها روی ۵۰۰۰ جفت‌سوال Quora}
\label{tab:grid}
\footnotesize
\begin{tabular}{lcccccccc}
\toprule
پیکربندی & روش شباهت‌سنجی & آستانه & دقت & بازیابی & F1 & صحت & کاهش مقایسه & زمان (ثانیه) \\
\midrule
\textbf{Configuration A} & MinHash+LSH & ۰٫۱۵ & ۰٫۵۴۴۸ & ۰٫۳۲۷۶ & ۰٫۴۰۹۲ & ۰٫۶۳۸۴ & ٪۷۰٫۶۲ & ۳٫۸۲ \\
(پیکربندی بهینه) & SimHash & ۱۸ & ۰٫۵۵۲۹ & ۰٫۵۵۵۲ & ۰٫۵۵۴۰ & ۰٫۶۵۸۴ & ٪۰٫۰۰ & ۱٫۰۴ \\
\midrule
\textbf{Configuration B} & MinHash+LSH & ۰٫۱۵ & ۰٫۵۴۶۲ & ۰٫۳۳۴۴ & ۰٫۴۱۴۸ & ۰٫۶۴۱۰ & ٪۶۶٫۶۰ & ۴٫۲۱ \\
& SimHash & ۱۶ & ۰٫۵۶۹۳ & ۰٫۴۴۴۸ & ۰٫۴۹۹۴ & ۰٫۶۵۵۰ & ٪۰٫۰۰ & ۱٫۰۱ \\
\midrule
\textbf{Configuration C} & MinHash+LSH & ۰٫۲۰ & ۰٫۵۳۲۴ & ۰٫۲۸۳۶ & ۰٫۳۷۰۱ & ۰٫۶۳۴۲ & ٪۶۶٫۶۰ & ۴٫۱۵ \\
& SimHash & ۱۴ & ۰٫۵۸۰۳ & ۰٫۳۳۲۸ & ۰٫۴۲۳۰ & ۰٫۶۴۹۶ & ٪۰٫۰۰ & ۰٫۹۸ \\
\midrule
\textbf{Configuration D} & MinHash+LSH & ۰٫۲۵ & ۰٫۵۰۶۸ & ۰٫۱۷۶۳ & ۰٫۲۶۱۶ & ۰٫۶۲۵۰ & ٪۸۶٫۱۸ & ۳٫۹۵ \\
(پایه اولیه) & SimHash & ۲۵ & ۰٫۴۷۰۴ & ۰٫۸۸۲۳ & ۰٫۶۱۳۶ & ۰٫۵۸۴۸ & ٪۰٫۰۰ & ۱٫۰۵ \\
\bottomrule
\end{tabular}
\end{table}

\begin{takeawaybox}
پیکربندی بهینه انتخاب‌شده (\lr{Configuration A}) توانست نرخ بازیابی MinHash را تقریباً دو برابر کند (از ٪۱۷٫۶۳ به ٪۳۲٫۷۶) و هم‌زمان دقت SimHash را به میزان قابل توجهی ارتقا داده و تعادل بهتری در هر دو مسیر برقرار کند.
\end{takeawaybox}

\section{تحلیل خطا و تفسیر عمیق موارد شکست}

\subsection{خطای مثبت کاذب (False Positives)}
مثبت کاذب زمانی رخ می‌دهد که دو جمله واژگان مشترک زیادی دارند اما معنای متفاوتی را می‌رسانند.
\begin{itemize}
\item \textbf{سوال اول}: \lr{"How can I make money online?"}
\item \textbf{سوال دوم}: \lr{"How can I not make money online?"}
\end{itemize}
تفاوت این دو جمله تنها در کلمه نفی \lr{not} است. از آنجا که \lr{SimHash} در سطح تک‌واژه فرکانسی عمل می‌کند، توزیع علامت‌های بردار این دو جمله بسیار نزدیک به هم شده و فاصله همینگ ۱۱ را تولید می‌کند که منجر به مثبت کاذب می‌گردد.

\subsection{خطای منفی کاذب (False Negatives)}
منفی کاذب زمانی رخ می‌دهد که بازنویسی معنایی عمیق بدون هم‌پوشانی لغوی انجام شود.
\begin{itemize}
\item \textbf{سوال اول}: \lr{"What is the best way to lose weight?"}
\item \textbf{سوال دوم}: \lr{"How can I reduce body fat?"}
\end{itemize}
با وجود تشابه مفهومی کامل، هیچ شینگل مشترک واژگانی بین دو جمله وجود ندارد که منجر به شباهت جاکارد صفر و شناسایی نشدن توسط مسیر MinHash می‌گردد.

\section{تجربیات مهندسی، چالش‌ها و درس‌های توسعه}

در مسیر توسعه این موتور تشخیص مشابهت اسناد، یادگیری‌های تجربی مهمی حاصل شد که در قالب چالش‌ها و مسیرهای بهینه‌سازی به شرح زیر دسته‌بندی می‌شوند:

\subsection{چالش فیلترینگ شدید ایست‌واژه‌ها (تجربه منفی)}
در مراحل اولیه پیاده‌سازی خط لوله پیش‌پردازش، از یک لیست ایست‌واژه سنتی و گسترده استفاده شد. این فیلترینگ شدید منجر به حذف توکن‌های پرسش و نفی (مانند \lr{what}، \lr{how}، \lr{why} و \lr{not}) گردید. این اقدام به ظاهر درست، ساختار جملات کوتاه دو جفت‌متن را متقارن کرده و دقت مسیر \lr{SimHash} را به حدود ٪۴۷٫۰۴ کاهش داد که حاصل آن تولید ۱۸۵۴ مورد مثبت کاذب (False Positives) روی ۵۰۰۰ جفت‌سوال Quora بود. این تجربه نشان داد که در متون کوتاه و سوالات مفهومی، کلمات ربط و نفی نقش ساختاری در متمایزسازی معنا دارند.

\subsection{مسیر بهبود و بهینه‌سازی چندجانبه (راهکار و مسیر درست)}
پس از شناسایی ضعف فوق، خط لوله پیش‌پردازش به فرم محافظه‌کارانه بازنویسی شد تا توکن‌های معناساز حفظ شوند. علاوه بر آن، با بهینه‌سازی شبکه پارامترها و کاهش آستانه فاصله همینگ \lr{SimHash} به ۱۸، اثر نویز کلمات مشترک خنثی شد که این امر دقت را به ٪۵۵٫۲۹ ارتقا داد. همچنین با کاهش آستانه جاکارد مسیر \lr{MinHash+LSH} به $0.15$، کارکرد فیلترهای LSH در هم‌پوشانی‌های کوچک تثبیت شد و بازیابی به میزان ۲ برابر بهبود یافت.

\subsection{تحلیل کارایی کدهای اصلی در مقابل قابلیت‌های امتیازی}
جدول مقایسه نهایی بخش بعدی به وضوح نشان می‌دهد که استفاده از امضاهای ترکیبی کاراکتری در \lr{Hybrid SimHash} با ثبت الگوهای زیرکلمه‌ای، بازیابی را تا ٪۶۲٫۴۸ بالا می‌برد که تاییدی بر غلبه بر غلط‌های املایی و تفاوت‌های نگارشی است.

\section{بخش امتیازی: قابلیت‌های توسعه‌یافته پیشرفته (Bonus)}

سه ایده امتیازی از صفر برای ارتقای کارایی سیستم طراحی و پیاده‌سازی شدند:
\begin{enumerate}
\item \textbf{بهینه‌ساز پویای باندها (\texttt{AdaptiveLSH})}: الگوریتم خودکار ریاضی که با دریافت آستانه دلخواه، به صورت خودکار مقادیر بهینه $b$ و $r$ را بر اساس تقسیم‌پذیری طول امضا محاسبه می‌کند.
\item \textbf{بن‌واژه‌ساز فارسی قاعده محور (\texttt{persian\_lemmatize})}: ابزار سبک و بدون وابستگی که پسوندهای جمع (`ها`, `ان`, `ات`) و پیشوند استمراری فعلی (`می`) را حذف کرده تا هم‌پوشانی متون فارسی بهبود یابد.
\item \textbf{اثرانگشت ترکیبی کاراکتری (\lr{Hybrid SimHash})}: با ادغام ۳-گرام کاراکتری و کلمات، اثرانگشت‌ها را در برابر غلط‌های املایی یا تغییر حروف بسیار مقاوم می‌سازد.
\end{enumerate}

جدول \ref{tab:bonus} مقایسه عددی تاثیر قابلیت‌های امتیازی را روی ۵۰۰۰ جفت‌سوال نشان می‌دهد:

\begin{table}[H]
\centering
\caption{ارزیابی تاثیر الگوریتم‌های امتیازی (Bonus Configurations)}
\label{tab:bonus}
\begin{tabular}{lrrr}
\toprule
روش ارزیابی & دقت & بازیابی & F1-score \\
\midrule
\lr{Static MinHash+LSH} (Baseline) & ۰٫۵۴۴۸ & ۰٫۳۲۷۶ & ۰٫۴۰۹۲ \\
\lr{Adaptive MinHash+LSH} (Bonus) & ۰٫۵۴۴۸ & ۰٫۳۲۷۶ & ۰٫۴۰۹۲ \\
\lr{Standard SimHash} (Baseline) & ۰٫۵۵۲۹ & ۰٫۵۵۵۲ & ۰٫۵۵۴۰ \\
\lr{Hybrid SimHash} (Bonus) & \textbf{۰٫۵۴۷۲} & \textbf{۰٫۶۲۴۸} & \textbf{۰٫۵۸۳۴} \\
\bottomrule
\end{tabular}
\end{table}

بکارگیری امضاهای ترکیبی کاراکتری در \lr{Hybrid SimHash} منجر به افزایش بازیابی به \textbf{٪۶۲٫۴۸} (افزایش ۷ درصدی) و ارتقای اف‌وان به \textbf{٪۵۸٫۳۴} شد.

\section{جمع‌بندی}
در این مستند، موتور تحلیل تشابه اسناد به طور کامل ارزیابی و مستند گردید. با بهینه‌سازی پارامترهای ریاضی و استفاده از اثرانگشت ترکیبی کاراکتری، بازیابی و دقت سامانه در برابر بازنویسی‌های پیچیده بهبود یافت.

\begin{thebibliography}{9}
\bibitem{guide} فاطمه شاکری، «راهنمای پروژه سوم درس داده‌کاوی: موتور تشخیص سرقت ادبی و اسناد مشابه»، دانشگاه صنعتی امیرکبیر، بهار ۱۴۰۵.
\bibitem{lshnotes} فاطمه شاکری، «جزوه الگوریتم \lr{Locality Sensitive Hashing}»، دانشگاه صنعتی امیرکبیر.
\bibitem{textmining} فاطمه شاکری، «جزوه متن‌کاوی»، دانشگاه صنعتی امیرکبیر.
\end{thebibliography}

\end{document}
